% A metodologia deste TCC é dividida em metodologia de pesquisa e metodologia de desenvolvimento, que serão descritas brevemente a seguir.

% \subsection{Metodologia de Pesquisa}

% Conforme descrito por \citeonline{zanella2006metodologia}, a pesquisa apresenta diversos tipos relacionado à sua natureza, objetivos, abordagem e procedimentos. Ainda de acordo com o material, é possível dizer que a presente pesquisa possui:

% \begin{itemize}

% \item \textbf{Natureza:} aplicada, por ter como objetivo a solução de um problema humano;
% \item \textbf{Objetivos:} exploratória, em razão de ser uma pesquisa que busca ampliar o conhecimento sobre fenômenos da comunidade acadêmica, atrás de conhecer suas necessidades e buscar tratá-las; 
% \item \textbf{Abordagem:} quantitativa, visto que caracteriza-se pela ampla utilização de instrumental estatístico fundamental na análise de dados coletados;
% \item \textbf{Procedimentos:} procedimento bibliográfico para investigação do referencial teórico aliado a produção tecnológica no contexto de desenvolvimento de software.

% \end{itemize}

% \subsection{Metodologia de Desenvolvimento}

% A metodologia de desenvolvimento deste trabalho segue princípios de práticas ágeis e une características do Extreme Programming (XP) ao Scrum, Kanban e Lean Inception. 

% No que diz respeito a metodologias ágeis, os atributos de programação em pares, aplicação de testes automatizados com TDD e a refatoração foram aproveitados da metodologia XP \cite{highsmith_2000}. Já as características extraídas da metodologia Scrum foram os conceitos de sprints iterativas e incrementais, junto das reuniões diárias \cite{schwaber_scrum_2013}. Vale destacar que esta última metodologia coube adaptações no contexto de uma equipe reduzida. 

% O Kanban, por sua vez, foi adotado como uma ferramenta de gestão visual para acompanhar o andamento das tarefas. Por fim, o Lean Inception estruturou a definição e priorização do Mínimo Produto Viável (MVP), que possibilitou a conformidade entre os envolvidos sobre a ideia do produto final.

% Maiores detalhes sobre a metodologia de desenvolvimento serão fornecidas no Capítulo 3 deste projeto.